\subsubsection{Online admission, direct revaluation, and eviction}
\label{sec:placement-plan}

\paragraph{Compare current plans, not historical credits.}
The Master re-evaluates placement at the events described in
Section~\ref{sec:system-overview}. Each decision uses one versioned snapshot
of the DAG, demand summaries, committed copies, queues, leases, and byte
ledgers. A cached estimate is reusable only while these dependencies are
unchanged. There is no rent level, accumulated credit, or renewal rule.
Tool waiting alone neither deletes history nor pins a particular tier;
normal reuse and failure exposure both remain in Equation~\ref{eq:value}.

Direct revaluation is a design choice, not a claim of optimal search.
The isolated comparison in Section~\ref{sec:eval-rental} retains the
rejected rental alternative: historical credit screening increased
modeled miss-plus-movement cost without reducing replacement-value checks.
That additive example does not validate the complete DAG controller or
measure its CPU cost. The runtime therefore bounds candidate and victim
evaluations and reports both decision work and resulting placement regret.

\paragraph{Construct a feasible replacement.}
Each candidate is a legal promotion or newly sealed checkpoint group,
including the dependencies needed to advertise its new frontier. The
safe victim set $Q$ excludes replicas with outstanding physical read
leases or copy dependencies. Demotion needs a committed lower copy or a
feasible copy-and-commit schedule; last-copy deletion needs an explicit
valid fallback. Logical ownership by a running request does not
permanently pin a physical location. Shared bytes are charged once, while
the plan's effects on all owners are evaluated together.

For each candidate, the Master constructs a private victim plan. It scans
a bounded, rotating set of eligible groups in the affected tier, ordered
by immutable object ID from a persistent cursor. Starting from the current
placement, it computes each feasible removal's additional cost through
Equation~\ref{eq:value} per byte released in that tier. It appends the
least-loss removal, breaking ties by fewer moved bytes and then object ID,
and re-evaluates subsequent removals against the changed tentative plan.
This continues until sufficient room exists for the incoming group and
its transfer workspace, or the configured work budget is exhausted.
The cursor resumes at the next eligible ID on later decisions; a bounded
scan can miss better victims and supplies no global optimality guarantee.

A DRAM demotion that lacks storage space first constructs a lower-tier
relief plan by the same rule, with storage deletion as the final tier.
It includes those victims, fallback changes, and copy dependencies in
the same plan. It never assumes an infinite lower tier. A tentative
removal is eligible only if its transition can be scheduled within the
physical and reserved byte limits; incoming copies wait for real release.
If no such plan fits the work budget, this candidate is skipped. The
ratio only constructs the victim set: the final score is the complete
candidate-plus-victim cost, not the sum of independent block scores.

\begin{figure}[t]
 \centering
 \includegraphics[width=\columnwidth]{figures/placement-competition.pdf}
 \caption{Direct replacement comparison in a constructed additive
 example with equal-size groups and cost units of seconds. Promoting x
 saves 6; demoting y or z costs 2 or 5; either transition costs 1.
 Thus the complete changes are $1+2-6=-3$ and $1+5-6=0$.
 A leased replica is excluded even if its removal cost is lower.
 With a margin below 3, the first plan is eligible for reservation;
 physical release still waits for commit and reader drainage. General
 common-failure plans use the joint DAG, not this additive shortcut.}
 \label{fig:placement-competition}
 \Description{Current values compare two complete candidate and victim
 plans; a leased object is excluded and only a profitable safe plan proceeds.}
\end{figure}

\paragraph{Choose and apply one plan.}
Among the fully evaluated plans, the Master selects the smallest
$\Delta J$, breaking ties by fewer moved bytes and stable object IDs.
An optional plan must improve predicted cost by more than the configured
uncertainty allowance plus migration hysteresis. These allowances are
heuristic guards unless separately calibrated; they are not statistical
confidence bounds. Rejected plans change neither placement nor capacity.
Before acceptance, the Master rechecks snapshot versions and reserves
the immediately executable destinations and transfer workspace. It queues
dependent actions in copy--commit--redirect--drain--release order and
starts incoming copies only when their space is physically available.
Later completions trigger a new snapshot; tentative plans from an older
placement are not reused as if nothing had changed.

Required capacity relief has no incoming bidder. The same current-loss
rule constructs a safe relief plan until the target is met, re-evaluating
after each tentative removal. It may accept a positive cost to satisfy
the capacity constraint. Storage deletion considers unowned or superseded
state first. Further relief may withdraw an active checkpoint claim,
but must first publish a complete older manifest or a prompt-and-token-log
fallback that does not depend on the deleted object. All owners' costs
enter the plan, and old-generation readers must drain before deletion.
If no executable plan exists within the budget, new admission pauses;
the old frontier remains valid and a later event resumes the search.

\par\medskip\noindent\begin{minipage}{\columnwidth}
\raggedright
 \refstepcounter{designalgorithm}\label{alg:placement}
 \textbf{Algorithm \thedesignalgorithm: Direct placement revaluation}
 \par\smallskip
 \begin{tabularx}{\columnwidth}{@{}r@{\quad}X@{}}
 \toprule
 1 & On a review event, snapshot versions, copies, queues and byte ledgers \\
 2 & Refresh causal demand and affected placement values; reuse existing objects \\
 3 & If capacity must shrink, construct a safe current-loss relief plan \\
 4 & Otherwise, form bounded legal sparse/layer candidates; best $\leftarrow$ none \\
 5 & For each candidate within the work budget: \\
 6 & \quad Construct safe victims by re-evaluated loss per released byte \\
 7 & \quad Include lower-tier victims and fallback changes; skip if infeasible \\
 8 & \quad Evaluate the complete plan with Equation~\ref{eq:value} \\
 9 & \quad Keep the lowest-cost plan that clears uncertainty and hysteresis \\
 10 & If no plan qualifies, retain placement; on pressure, pause new admission \\
 11 & Recheck versions; on mismatch, reschedule from a fresh snapshot \\
 12 & Reserve executable space and enqueue dependent transitions \\
 13 & On completion/failure, reconcile actual state and trigger a fresh review \\
 \bottomrule
 \end{tabularx}
 \par\smallskip\noindent\textit{Predicted release is not free memory; a plan is not a committed copy.}
 \end{minipage}\par\medskip

Profiles carry model/layout, size, tail, topology, queue load, support
range, uncertainty, and expiry. Missing or stale profiles freeze optional
promotions and use bounded LRU for required relief under the same safety
protocol. Candidate construction, each victim comparison, transfer
profiling, and complete DAG passes count against the work budget; report
their counts, decision latency, and the loss from bounded search. The
isolated model's replacement checks are not measurements of these costs.

